\section{Independent specification and history bridges}
\label{sec:history-bridge}

Representation correctness answers whether a stored state is explained by the
issued event set. The independent specification answers whether that explanation
is a permitted sequential history. These are separate obligations. In particular,
convergence of implementation replays of the same history does not establish that
any of their query answers belongs to the intended datatype.

Use the independent history language and projection of Definition~\ref{def:history}.
Write $S$ for that language, retain its projection $f$, and abbreviate
\[
 F(\pi)=u^*(\sigma_0,\pi).
\]
This expression contains no merges. Write $L_{C,E}$ for the implementation order
chosen for context $C$ and endpoint event set $E$: the raw order of
Definition~\ref{def:order}, or its separately justified invariant-scoped
replacement. Its separate specification-visibility obligation
uses $\operatorname{sv}^{f,S}_C$ from Definition~\ref{def:history}. Event identities remain
in the witness even when the independent specification ignores their metadata.

\subsection{A sufficient global route}

A local simulation into an independent deterministic sequential machine $M$
for payload labels consists of a relation $R_{\mathrm{seq}}$, initial machine state $m_0$,
machine update $w$ and query $\operatorname{read}_M$, with these properties:
\begin{align*}
 &R_{\mathrm{seq}}(\sigma_0,m_0),\\
 &R_{\mathrm{seq}}(s,m)\Longrightarrow
   R_{\mathrm{seq}}(u(s,e),w(m,e.o))
   &&\text{for every event }e,\\
 &R_{\mathrm{seq}}(s,m)\Longrightarrow
   \operatorname{read}(s,q)=\operatorname{read}_M(m,q)
   &&\text{for every query }q.
\end{align*}
There is no issuance, freshness or invariant premise on the update clause.
For $f(e)=e.o$, induction establishes the unrestricted sequential-history premise
\begin{equation}
 \forall\pi,q,\qquad
 \operatorname{labels}_f(\pi)\cdot
 [\operatorname{qry}\bigl(q,\operatorname{read}(F(\pi),q)\bigr)]\in S.
 \label{eq:total-fold-sound}
\end{equation}
% SequentialSimulation.lean: SequentialSimulation; SequentialSimulation.sound.
% RALinearizability.lean: FoldHistorySound; uniform_ra_of_replay_total;
% ra_of_canonical_total; certifiedRA_of_join_total.
Together with a stored replay witness and the replay laws, this suffices for the
literal criterion: choose the replay word and substitute its equality with the
stored state. The unrestricted premise concerns \emph{merge-free folds}; it does
not assert that every permutation respects $L_{C,E}$ or is an admissible
linearization.

For the stronger criterion, the global bridge additionally records
commutation compatibility:
\begin{equation}
 \operatorname{Commutes}_{\mathrm{impl}}(a,b)
 \Longrightarrow f(a)\bowtie_S f(b).
 \label{eq:commutation-compatibility}
\end{equation}
Here implementation commutation must be the commutation notion used in the
chosen order. By contraposition, every specification-conflicting visibility
edge is then an implementation-conflicting visibility edge. The same replay
word respects both relations. Equation~\eqref{eq:total-fold-sound} still supplies
history admission; compatibility alone cannot supply it.
% CommutationBridge.lean: CommutationCompatibility;
% projectedSpecVisibility_sub_paperOrder; eventSpecificationRA_of_compatibility.
% GuardedHistoryBridge.lean: of_canonical; ScopedCertificate.ofTotal.

\begin{example}[A convergent set with a no-op remove]\label{ex:broken-set}
Consider a one-element set implemented by a finite set of addition timestamps, initially empty.
An add inserts its timestamp; a remove leaves the state unchanged; the query
returns whether the timestamp set is nonempty. Retain the ordinary observed-remove
three-way merge
\[
 m(l,a,b)=(l\cap a\cap b)\cup(a\setminus l)\cup(b\setminus l).
\] Let the independent sequential machine instead implement add
as setting presence to true and remove as setting it to false.

The real, single-replica execution is
\[
 \operatorname{add}\ ;\ \operatorname{remove}\ ;\
 \operatorname{query}(\mathrm{true}).
\]
All implementation updates commute, so with an empty arbitration policy the
literal implementation order is empty. The witness
$[\operatorname{remove},\operatorname{add}]$ explains the wrong answer in the
independent language and satisfies the literal criterion. Indeed, this mutant
converges and satisfies the literal criterion on every raw ordinary execution;
it does not satisfy unrestricted fold soundness.

In the independent set language, add and remove do not contextually commute.
Their causal order is therefore required by $\operatorname{sv}^{f,S}_C$. The only respecting
witness is add before remove, whose query answer is false. The stronger
criterion rejects the displayed execution. This is a certified counterexample
to identifying literal RA with preservation of sequential specification
conflicts; it is not a counterexample to a theorem that assumes
Equation~\eqref{eq:total-fold-sound}.
\end{example}
% CriterionCounterexample.lean: all_commute; mutant_join; mutant_rawRA;
% wrong_answer_ra_linearizable; causal_history_rejects_true;
% wrong_answer_not_specification_ra; criteria_differ_on_reachable_mutation.

\subsection{Choosing a permitted history within a certified execution}

Global fold soundness can be unsuitable for a datatype whose public operations
have guards: an arbitrary word may name a nonexistent birth, reuse an identity,
or remove a live non-head. Restricting the implementation's representation
invariant does not automatically prove the independent language's guards.
A second route proves existence of a suitable history on the supported,
causally closed event sets of certified executions.

\begin{definition}[Execution-scoped history adequacy]\label{def:history-adequacy}
Fix an issuance discipline $G$, an implementation order $L$, and an independent
language $S$. They satisfy execution-scoped history adequacy when, for every
$G$-certified configuration $C$, every stored version $(s,E)$ and every query $q$,
there is a word $\pi$ such that
\begin{align*}
 &\pi\text{ enumerates }E\text{ exactly once},\\
 &\pi\text{ respects }L_{C,E}\text{ and }\operatorname{sv}^{f,S}_C,\\
 &\operatorname{labels}_f(\pi)\cdot
   [\operatorname{qry}\bigl(q,\operatorname{read}(F(\pi),q)\bigr)]\in S.
\end{align*}
The quantifier order is $\forall q\,\exists\pi$. The accepted answer describes the
merge-free fold, rather than the stored state $s$. Both order obligations and
language admission use the same $\pi$.
\end{definition}
% ScopedHistoryBridge.lean: AbstractMRDT.ExecutionHistoryAdequacy.
% RALinearizability.lean: HistoryAdequate; IssuedHistoryAdequacy (literal variant).
% SpecificationVisibility.lean: IssuedSpecificationHistoryAdequacy.

\begin{theorem}[Representation and history bridge]\label{thm:history-bridge}
Fix the datatype, issuance discipline, independent language, and chosen
implementation replay order. Define exact concrete canonicality by
\[
 \operatorname{Can}_{C,E}(s)\quad\Longleftrightarrow\quad
 \exists\pi.\ \pi\text{ enumerates }E\ \land\
 \pi\text{ respects }L_{C,E}\ \land\ F(\pi)=s.
\]
The concrete specialization of the history bridge has these assumptions:
\begin{enumerate}[nosep]
 \item the chosen replay order has a proved uniqueness property: any two
       exact canonical states for the same supported event set agree on all queries;
 \item every stored event set of the certified execution is supported, and its
       stored state is canonical for that set;
 \item execution-scoped history adequacy of
       Definition~\ref{def:history-adequacy} holds.
\end{enumerate}
Then every stored version satisfies the criterion with independent
specification visibility: one word enumerates its original events, respects
both required relations, and explains its stored query answer in $S$.
Consequently all active heads satisfy that criterion. If the assumptions are
established at every configuration of a finite certified execution, the result
holds throughout that execution, including executions with recursively
synthesized virtual merge bases.
\end{theorem}
% ScopedHistoryBridge.lean: AbstractMRDT.versions_of_scoped_history;
% ScopedVCConditions.versionsV; ScopedVCConditions.executions;
% ScopedVCConditions.executionsV.
% GuardedHistoryBridge.lean: Guarded.versions_of_scoped_history;
% Guarded.ScopedCertificate.versions; Guarded.ScopedCertificate.executionsV.

\noindent\emph{Proof.}
Adequacy supplies a word $\pi$ whose fold is canonical by its permutation and
order properties. Replay uniqueness identifies the queries of that fold and
the separately canonical stored state. Substitute this query equality in the
accepted history. Its word, permutation and two order properties do not change.
Apply the argument at each visited configuration for the execution claim.
The stored-state premise is where the merge argument enters; history adequacy
alone cannot discharge it.\hfill$\square$

This uniqueness assumption must be established for the chosen implementation
order; it is not automatic from scoped laws. For example,
Theorem~\ref{thm:scoped-replay} requires causally ready implementation replays;
it does not give equality of every pair of permutations respecting only the
invariant order. A concrete port may instead prove an exact representation
relation between its selected independent history and the stored live state,
as in Example~\ref{ex:queue}.

\begin{example}[Anchored enqueue with an independent FIFO]\label{ex:queue}
Enqueue records its issuer's observed tail and dequeue its observed head. The
implementation keeps ordered live records and carried coordinates; it stores
no removed birth records. The independent FIFO appends a fresh tagged value,
ignoring its recorded tail; it removes a named head, permits repeated removal
of an already removed identity, and rejects removal of a live non-head.
The absent-removal case requires a prior birth in the sequential prefix;
invented identities cannot be removed. Lack of tombstones does not mean lack
of metadata: carried coordinates preserve placement information and can grow
with their ancestry.

Head-only issuance implies that removed births form a causal downset among
enqueue events. Order the original dequeue events chronologically. Immediately
before an identity's first dequeue, emit its birth; emit subsequent duplicate
dequeues without another birth. Append the surviving births in their stored
coordinate order. The FIFO is empty between the removal blocks. The downset
property preserves causal enqueue conflicts, and surviving coordinates preserve
the remaining causal enqueue order. This one word has the exact final tagged
live list and satisfies both order obligations. The resulting theorem applies
to certified ordinary and recursive virtual executions with the unchanged
public tagged-head query.
\end{example}
% AnchoredQueueCausality.lean: removed_births_causal_downset.
% AnchoredQueueRanks.lean: DeadEndpoint.causalRanks; exists_deadEndpoint.
% AnchoredQueueSchedule.lean: deadWord_legal; schedule_legal; schedule_fold.
% AnchoredQueueCertificate.lean: History.stored_publicWitness;
% History.correct; History.correctV; History.executions; History.executionsV.

\begin{example}[Fugue and a live-anchor specification]\label{ex:fugue}
The original Fugue issuer permits insertion naming a deleted birth. A certified
trace creates a birth, deletes it, then inserts a replacement anchored to that
birth. Its independent live-anchor language permits the replacement before the
deletion, but rejects it after the deletion. Specification visibility requires
the actual causal deletion-before-replacement order. No admitted word satisfies
that requirement, for any representation-state domain or arbitration policy.
The relevant implementation effectors nevertheless commute on the proved closed
invariant domain. Thus representation invariants repair an implementation-order
issue without repairing this issuance/specification mismatch. A successful
claim must reconcile those two interfaces explicitly.
\end{example}
% CertifiedFugueInvariantObstruction.lean: certified_failure;
% implementation_spec_separation; valid_execution_counterexample.
